\abstract % Структурный элемент: РЕФЕРАТ

\keywords{УПРАВЛЕНИЕ ИНФРАСТРУКТУРОЙ, ARCHITECTURE-AS-CODE, DESIRED STATE, ACTUAL STATE, DRIFT DETECTION, RECONCILE, CONTROL LOOP, INVENTORY, SRE, DEVOPS}

Выпускная квалификационная работа посвящена автоматизированному управлению инфраструктурой. Актуальность темы обусловлена необходимостью контроля соответствия между архитектурным описанием и фактическим состоянием без участия человека. В распределенных системах такая согласованность нарушается из-за ручных изменений, неполных процедур развертывания, сбоев и различий между средами. Возникающие расхождения рассматриваются как drift и требуют своевременного обнаружения и устранения.

Цель работы --- спроектировать и реализовать автоматизированную платформу управления инфраструктурой с поддержкой подхода Architecture-as-Code, обеспечивающую формирование desired state, получение actual state, обнаружение drift и инициирование reconcile для приведения инфраструктуры к заданному состоянию. Для этого выполнен анализ предметной области, проведено сравнение существующих решений, сформулированы требования к платформе, спроектирована архитектура и реализованы основные компоненты программного решения.

Практическим результатом работы стала MVP-реализация Infrastructure Management Platform, поддерживающая контроль наличия компонент \texttt{node\_exporter} и \texttt{cadvisor}. На данном сценарии подтверждена работоспособность полного цикла управления: архитектурное описание используется как источник desired state, inventory-слой формирует actual state, drift detection выявляет отсутствие требуемого компонента, а reconcile-контур инициирует его восстановление.

Предложенная платформа может рассматриваться как основа для дальнейшего развития внутреннего control plane, в котором Architecture-as-Code используется не только как способ документирования, но и как источник данных для автоматического управления инфраструктурой.
